\begin{figure}[tbp]
\centering
\begin{tikzpicture}
\begin{axis}[width=.95\textwidth,height=7cm,ymin=0,ymax=100,ylabel={Development F1 (\%)},
 xtick={0,1,2,3},xticklabels={Static,Threshold 144,Refit 144/144,Refit 288/240},
 x tick label style={rotate=15,anchor=east,font=\small},
 legend style={at={(.5,1.02)},anchor=south,legend columns=2,column sep=6pt},
 ymajorgrids=true,grid style={gray!20}]
\addplot[ChromecastBlue,thick,mark=*] table[x=mode,y=f1,col sep=comma]{generated/legacy_UNSW-Chromecast.csv};
\addplot[SamsungOrange,thick,dashed,mark=square*] table[x=mode,y=f1,col sep=comma]{generated/legacy_UNSW-Samsung-Camera.csv};
\legend{Chromecast,Samsung camera}
\end{axis}
\end{tikzpicture}
\caption{Historical DM-only development comparisons. Refit labels give attempt
cadence/minimum accepted windows; every adaptive variant has a 288-window buffer.
The lines connect distinct configurations, not a temporal learning trajectory.
These development-selected runs differ from the conservative closing replay.}
\label{fig:legacy}
\end{figure}
